% Clase 7 — el problema resuelto: esfera conductora descargada de radio a en un
% campo uniforme. Se marcan las DOS condiciones de borde y el dato extra (Q=0)
% que hizo falta para cerrar la constante.
% Las dos CB van SIMETRICAS a media altura y con relleno blanco: apiladas arriba
% se comen el arco de theta y el rotulo de r (pasa en el render, no en el codigo).
\begin{tikzpicture}[scale=1]

  \draw[figaxis] (0,-2.5) -- (0,3.4);
  \node[figlblaux] at (0,3.62) {$z$};

  \foreach \xx in {-3.5,-2.85,2.85,3.5}{
    \draw[figvecaux] (\xx,-2.3) -- (\xx,2.9);
  }
  \node[figlblaux] at (-3.5,3.22) {$\vec E_0$};

  % esfera conductora
  \fill[figgray!25] (0,0) circle (1.0);
  \draw[figline, line width=1.1pt] (0,0) circle (1.0);
  \node[figlblaux] at (0,-0.36) {$\vec E=0$};
  \draw[figvecaux] (0,0) -- (218:1.0);
  \node[figlblaux, fill=figgray!25, inner sep=1.2pt] at (-0.4,-0.06) {$a$};
  \foreach \ang in {62,76,90,104,118}{\node[figpos] at (\ang:1.0) {};}
  \foreach \ang in {242,256,270,284,298}{\node[figneg] at (\ang:1.0) {};}

  % punto de observacion, bien arriba, con r y theta libres de las CB
  \coordinate (P) at (1.15,2.55);
  \node[figneutra] at (P) {};
  \draw[figvec] (0,0) -- (P);
  \node[figlbl, fill=white, inner sep=1.5pt] at (0.98,1.72) {$r$};
  \draw[figgray, line width=0.6pt] (0,1.95) arc (90:65.7:1.95);
  \node[figlblaux, fill=white, inner sep=1.2pt] at (0.3,2.12) {$\theta$};

  % las dos condiciones de borde, simetricas y a media altura
  \node[figlbl, figred, fill=white, inner sep=2pt] at (-2.55,0.9) {\textbf{CB 1}};
  \node[figlblaux, align=center, fill=white, inner sep=2pt] at (-2.55,0.28)
    {$\phi(a,\theta)=\phi_0$\\ (equipotencial)};
  \draw[figvecaux] (-1.78,0.5) -- (-0.88,0.42);

  \node[figlbl, figred, fill=white, inner sep=2pt] at (2.75,0.9) {\textbf{CB 2}};
  \node[figlblaux, align=center, fill=white, inner sep=2pt] at (2.75,0.28)
    {$\phi\to-E_0\,r\cos\theta$\\ cuando $r\to\infty$};

  % el dato extra
  \node[figlbl, align=center] at (0,-3.2)
    {\textbf{Dato extra que hace falta:} la esfera est\'a \emph{descargada}, $Q=0$.};
  \node[figlblaux, align=center] at (0,-3.95)
    {Sin \'el $\phi_0$ queda libre: cierran el problema las dos CB\\[0.15em]
     \emph{m\'as} el desarrollo multipolar (el t\'ermino en $1/r$ vale $Q/4\pi\varepsilon_0$).};

  \node[figlbl, align=center] at (0,-4.95)
    {$\phi(r,\theta)=-E_0\left(r-\dfrac{a^{3}}{r^{2}}\right)\cos\theta
     \qquad (r\ge a)$};

\end{tikzpicture}
